\documentclass[oneside,a4paper,14pt]{extarticle}
\usepackage[T1,T2A]{fontenc}
\usepackage[a4paper,letterpaper,top=20mm,bottom=20mm,left=30mm,right=10mm]{geometry}
\usepackage[utf8]{inputenc}
\usepackage[russian]{babel}
\usepackage{textcomp}
\usepackage{indentfirst}
\usepackage{graphicx}
\usepackage{mwe}
\usepackage{wrapfig}
\usepackage{caption}
\usepackage{amsmath}
\usepackage{amsfonts}
\usepackage{amsthm}
\usepackage[all]{xy}
\usepackage[breaklinks]{hyperref}
\usepackage{titlesec}
\titleformat{\section}{\normalsize\bfseries}{\thesection}{1em}{}
\titleformat{\subsection}{\normalsize\bfseries}{\thesubsection}{1em}{}
\titleformat{\subsubsection}{\normalsize\bfseries}{\thesubsection}{1em}{}
\renewcommand\baselinestretch{1.33}\normalsize
\setlength{\parindent}{1.25cm}
\usepackage{indentfirst}

\begin{document}
    \newpage\thispagestyle{empty}
    \begin{center}
        МИНИСТЕРСТВО НАУКИ И ВЫСШЕГО ОБРАЗОВАНИЯ\\
            РОССИЙСКОЙ ФЕДЕРАЦИИ
            ФЕДЕРАЛЬНОЕ ГОСУДАРСТВЕННОЕ БЮДЖЕТНОЕ\\
            ОБРАЗОВАТЕЛЬНОЕ
            УЧРЕЖДЕНИЕ ВЫСШЕГО ОБРАЗОВАНИЯ\\
            «ВЯТСКИЙ ГОСУДАРСТВЕННЫЙ УНИВЕРСИТЕТ»\\
            Институт математики и информационных систем\\
            Факультет автоматики и вычислительной техники\\
            Кафедра электронных вычислительных машин
    \end{center}
    \vspace{20mm}

    \begin{center}
        Отчёт по лабораторной работе №6\\
        по дисциплине\\
        <<Программирование>>\\
    \end{center}
    \vspace{48mm}

    Выполнил студент гр. ИВТб-1302-06-00 \hspace{10mm} \rule[-0,5mm]{30mm}{0.15mm}\,/Изместьев М.Н./


    Проверил доцент кафедры ЭВМ \hfill  \rule[-0,5mm]{30mm}{0.15mm}\,/Баташев П.А./

    \vfill
    \begin{center}
        Киров\\
        2026
    \end{center}

    \newpage\thispagestyle{empty}

\begin{flushleft}
\section*{Цель работы:}
\noindent Освоение методов предварительной обработки и анализа больших наборов данных с использованием библиотеки NumPy: загрузка структурированных данных, выявление и очистка аномалий, групповая нормализация, расчёт скользящих статистик, создание производных признаков, робастная фильтрация выбросов и проверка логической целостности данных.

\section*{Задание.}

1. Загрузить CSV-файл с данными своего варианта в структурированный массив NumPy с явным указанием dtype.

2. Выполнить разведочный анализ: подсчитать NaN, Inf, объём памяти.

3. Провести векторизованную фильтрацию и очистку аномалий.

4. Реализовать группировку по ключевому полю и z-нормализацию внутри каждой группы.

5. Вычислить скользящее среднее и разницу между соседними элементами.

6. Создать производные признаки на основе исходных полей.

7. Провести условную агрегацию по группам с сохранением в CSV.

8. Создать лаговые признаки (сдвиг на 1 позицию через \texttt{np.roll}).

9. Выполнить групповую робастную замену выбросов методом IQR.

10. Проверить логическую целостность и выполнить частотный анализ с сжатием редких категорий.

\textbf{Вариант 10} --- Спортивная аналитика.

\textbf{Поля данных:}
\begin{itemize}
\item \texttt{ts} (int32) --- UNIX-timestamp замера
\item \texttt{athlete\_id} (int16) --- ID спортсмена
\item \texttt{dist} (float32) --- дистанция тренировки, км
\item \texttt{pace} (float32) --- темп, мин/км
\item \texttt{cal} (float32) --- калории
\item \texttt{zone} (uint8) --- пульсовая зона (1--5)
\end{itemize}
\end{flushleft}

\newpage
\textbf{Полный исходный код (lab\_variant10\_numpy.py)}
\small
\begin{verbatim}
from __future__ import annotations
import argparse
from pathlib import Path
import numpy as np


def main(input_csv: str) -> None:
    # 1) Загрузка и подготовка типов
    raw = np.genfromtxt(input_csv, delimiter=",", names=True,
                        dtype=None, encoding=None)
    dtype = np.dtype([
        ("ts",         np.int32),
        ("athlete_id", np.int16),
        ("dist",       np.float32),
        ("pace",       np.float32),
        ("cal",        np.float32),
        ("zone",       np.uint8),
    ])
    data = np.empty(raw.shape[0], dtype=dtype)
    for name in data.dtype.names:
        data[name] = raw[name].astype(data.dtype[name], copy=False)

    n_rows    = data.shape[0]
    bytes_used = data.nbytes
    mb_used   = bytes_used / (1024.0 ** 2)
    print(f"Строк: {n_rows}")
    print(f"Память: {bytes_used} байт ({mb_used:.3f} Мбайт)")

    numeric_fields = ["dist", "pace", "cal"]
    nan_counts = {f: int(np.isnan(data[f]).sum())  for f in numeric_fields}
    inf_counts = {f: int(np.isinf(data[f]).sum())  for f in numeric_fields}
    total_bad  = sum(nan_counts.values()) + sum(inf_counts.values())
    bad_share  = total_bad / (n_rows * len(numeric_fields))
    print(f"NaN по полям: {nan_counts}")
    print(f"Inf по полям: {inf_counts}")
    if bad_share > 0.03:
        print("ПРЕДУПРЕЖДЕНИЕ: доля NaN/Inf > 3%")

    np.save("stage1_structured.npy", data)

    # 2) Векторизованная фильтрация и очистка
    anomaly_mask = (
        (data["dist"] < 0) | (data["pace"] < 0)
        | np.isnan(data["dist"]) | np.isnan(data["pace"]) | np.isnan(data["cal"])
        | np.isinf(data["dist"]) | np.isinf(data["pace"]) | np.isinf(data["cal"])
    )
    an_count = int(anomaly_mask.sum())
    print(f"Аномалий: {an_count} ({an_count / n_rows:.2%})")

    cleaned = data[~anomaly_mask].copy()
    cleaned["cal"] = np.where(cleaned["cal"] < 0, 0.0, cleaned["cal"])

    q10, q90 = np.nanpercentile(cleaned["pace"], [10, 90])
    cleaned["pace"] = np.clip(cleaned["pace"], q10, q90)

    # 3) Группировка + нормализация
    group_ids, group_counts = np.unique(cleaned["athlete_id"],
                                        return_counts=True)
    print(f"Групп: {group_ids.size}")
    print(f"Размеры групп: min={group_counts.min()}, max={group_counts.max()}")

    mean_dist = np.empty(group_ids.size, dtype=np.float32)
    std_dist  = np.empty(group_ids.size, dtype=np.float32)
    mean_pace = np.empty(group_ids.size, dtype=np.float32)
    min_pace  = np.empty(group_ids.size, dtype=np.float32)

    for i, gid in enumerate(group_ids):
        gm = cleaned["athlete_id"] == gid
        mean_dist[i] = np.nanmean(cleaned["dist"][gm])
        std_dist[i]  = np.nanstd (cleaned["dist"][gm])
        mean_pace[i] = np.nanmean(cleaned["pace"][gm])
        min_pace[i]  = np.nanmin (cleaned["pace"][gm])

    z_dist = np.empty(cleaned.shape[0], dtype=np.float32)
    for i, gid in enumerate(group_ids):
        gm = cleaned["athlete_id"] == gid
        z_dist[gm] = (cleaned["dist"][gm] - mean_dist[i]) / (std_dist[i] + 1e-8)

    norm_dtype = cleaned.dtype.descr + [("z_dist", np.float32)]
    norm_data  = np.empty(cleaned.shape[0], dtype=norm_dtype)
    for n in cleaned.dtype.names:
        norm_data[n] = cleaned[n]
    norm_data["z_dist"] = z_dist
    np.save("stage3_normalized.npy", norm_data)
\end{verbatim}

\newpage
\small
\begin{verbatim}
    # 4) Скользящее окно и разница
    ratio    = norm_data["cal"] / np.maximum(norm_data["dist"], 1e-8)
    k        = 30
    csum     = np.cumsum(np.insert(ratio, 0, 0.0))
    mov      = (csum[k:] - csum[:-k]) / k
    mov_full = np.pad(mov, (k - 1, 0), mode="edge")

    pace_diff      = np.diff(norm_data["pace"])
    pace_diff_full = np.insert(pace_diff, 0, 0.0).astype(np.float32)

    win_dtype = norm_data.dtype.descr + [
        ("cal_per_km_ma30", np.float32), ("pace_diff", np.float32)]
    win_data  = np.empty(norm_data.shape[0], dtype=win_dtype)
    for n in norm_data.dtype.names:
        win_data[n] = norm_data[n]
    win_data["cal_per_km_ma30"] = mov_full.astype(np.float32)
    win_data["pace_diff"]       = pace_diff_full

    # 5) Feature engineering
    speed_kmh   = win_data["dist"] / np.maximum(win_data["pace"] / 60.0, 1e-8)
    cal_per_min = win_data["cal"]  / np.maximum(win_data["pace"], 1e-8)

    speed_kmh   = np.where(np.isfinite(speed_kmh), speed_kmh, 0.0).astype(np.float32)
    cal_per_min = np.where(np.isfinite(cal_per_min), cal_per_min,
                           np.nanmedian(cal_per_min)).astype(np.float32)

    feat_dtype = win_data.dtype.descr + [
        ("speed_kmh", np.float32), ("cal_per_min", np.float32)]
    feat_data  = np.empty(win_data.shape[0], dtype=feat_dtype)
    for n in win_data.dtype.names:
        feat_data[n] = win_data[n]
    feat_data["speed_kmh"]   = speed_kmh
    feat_data["cal_per_min"] = cal_per_min

    # 6) Условная агрегация по группам
    # Условие: зона >= 4, дистанция > медианной, темп < медианного
    global_dist_median = np.nanmedian(feat_data["dist"])
    cond_mask = (
        (feat_data["zone"] >= 4)
        & (feat_data["dist"] > global_dist_median)
        & (feat_data["pace"] < np.nanmedian(feat_data["pace"]))
    )

    conditional_stats    = np.empty((group_ids.size, 3), dtype=np.float64)
    conditional_stats_p90 = np.empty((group_ids.size, 4), dtype=np.float64)
    for i, gid in enumerate(group_ids):
        gm   = (feat_data["athlete_id"] == gid) & cond_mask
        vals = feat_data["cal_per_min"][gm]
        if vals.size == 0:
            conditional_stats[i]    = [gid, np.nan, np.nan]
            conditional_stats_p90[i] = [gid, np.nan, np.nan, np.nan]
        else:
            m   = np.nanmean(vals)
            med = np.nanmedian(vals)
            p90 = np.nanpercentile(vals, 90)
            conditional_stats[i]    = [gid, m, med]
            conditional_stats_p90[i] = [gid, m, med, p90]

    np.savetxt("conditional_stats.csv", conditional_stats,
               delimiter=",",
               header="group_id,mean_conditional,median_conditional",
               comments="", fmt=["%d", "%.6f", "%.6f"])
    np.savetxt("conditional_stats_with_p90.csv", conditional_stats_p90,
               delimiter=",",
               header="group_id,mean_conditional,median_conditional,p90_conditional",
               comments="", fmt=["%d", "%.6f", "%.6f", "%.6f"])

    # 7) Лаговые признаки
    prev_pace      = np.roll(feat_data["pace"], 1)
    prev_pace[0]   = feat_data["pace"][0]
    pace_lag_diff  = feat_data["pace"] - prev_pace

    lag_dtype = feat_data.dtype.descr + [
        ("pace_lag1", np.float32), ("pace_delta_lag1", np.float32)]
    lag_data  = np.empty(feat_data.shape[0], dtype=lag_dtype)
    for n in feat_data.dtype.names:
        lag_data[n] = feat_data[n]
    lag_data["pace_lag1"]       = prev_pace.astype(np.float32)
    lag_data["pace_delta_lag1"] = pace_lag_diff.astype(np.float32)

    grew = np.mean(pace_lag_diff > 0)
    fell = np.mean(pace_lag_diff < 0)
    same = np.mean(pace_lag_diff == 0)
    print(f"Доля роста pace: {grew:.2%}, падения: {fell:.2%}, "
          f"без изменений: {same:.2%}")

    signs, sign_counts = np.unique(
        np.sign(pace_lag_diff).astype(np.int8), return_counts=True)
    print(f"Распределение sign: {dict(zip(signs.tolist(), sign_counts.tolist()))}")
\end{verbatim}

\newpage
\small
\begin{verbatim}
    # 8) Групповая робастная замена выбросов (IQR) по pace
    replaced_per_group  = []
    replaced_mask_total = np.zeros(lag_data.shape[0], dtype=bool)
    robust_pace         = lag_data["pace"].copy()

    for gid in group_ids:
        gm      = lag_data["athlete_id"] == gid
        vals    = robust_pace[gm]
        q1, q3  = np.nanpercentile(vals, [25, 75])
        iqr     = q3 - q1
        lo, hi  = q1 - 1.5 * iqr, q3 + 1.5 * iqr
        out     = gm & ((robust_pace < lo) | (robust_pace > hi))
        med     = np.nanmedian(vals)
        robust_pace[out] = med
        replaced_per_group.append((int(gid), int(out.sum())))
        replaced_mask_total |= out

    lag_data["pace"] = robust_pace
    replaced_share   = replaced_mask_total.mean()
    print(f"Доля замен после IQR: {replaced_share:.2%}")
    print(f"Замены по группам (первые 10): {replaced_per_group[:10]}")

    # 9) Проверка логической целостности
    # Зона 5 не должна идти с очень медленным темпом (> 8.5 мин/км),
    # зона 1 — с экстремально быстрым (< 3.2 мин/км).
    invalid_logic = (
        ((lag_data["zone"] == 5) & (lag_data["pace"] > 8.5))
        | ((lag_data["zone"] == 1) & (lag_data["pace"] < 3.2))
    )
    invalid_count = int(invalid_logic.sum())
    print(f"Логических нарушений: {invalid_count} "
          f"({invalid_count / lag_data.shape[0]:.2%})")

    lag_data["zone"] = np.where(
        (lag_data["pace"] < 4.0) & (lag_data["zone"] < 3),
        3, lag_data["zone"]).astype(np.uint8)
    lag_data["zone"] = np.where(
        (lag_data["pace"] > 8.0) & (lag_data["zone"] > 3),
        3, lag_data["zone"]).astype(np.uint8)

    # 10) Частотный анализ и сжатие редких категорий zone
    cats, cnts = np.unique(lag_data["zone"], return_counts=True)
    freq       = cnts / lag_data.shape[0]
    rare_cats  = cats[freq < 0.01]
    is_rare    = np.isin(lag_data["zone"], rare_cats)
    lag_data["zone"] = np.where(is_rare, 0, lag_data["zone"]).astype(np.uint8)
    print(f"Редкие категории zone: {rare_cats.tolist()}")

    np.save("final_preprocessed.npy", lag_data)
    print(
        "Готово. Сохранены файлы: stage1_structured.npy, "
        "stage3_normalized.npy, conditional_stats.csv, "
        "conditional_stats_with_p90.csv, final_preprocessed.npy"
    )


if __name__ == "__main__":
    parser = argparse.ArgumentParser()
    parser.add_argument("--input", default="data.csv")
    args = parser.parse_args()
    if not Path(args.input).exists():
        raise FileNotFoundError(f"Не найден файл {args.input}")
    main(args.input)
\end{verbatim}

\normalsize
\newpage
\begin{flushleft}
\section*{Описание этапов обработки}

\textbf{Этап 1. Загрузка и подготовка типов.} Данные читаются через \texttt{np.genfromtxt} с заголовками. Затем создаётся структурированный массив с жёсткими типами (int32, int16, float32, uint8) и выполняется перекастование полей. Подсчитываются NaN и Inf, вычисляется объём занятой памяти.

\textbf{Этап 2. Фильтрация аномалий.} Векторная маска объединяет все недопустимые условия: отрицательные dist или pace, NaN, Inf. Очищенный массив создаётся через булеву индексацию \texttt{data[\textasciitilde mask]}. Темп обрезается по перцентилям [10, 90] через \texttt{np.clip}.

\textbf{Этап 3. Группировка и z-нормализация.} \texttt{np.unique} возвращает уникальные \texttt{athlete\_id}. Для каждой группы вычисляются mean и std дистанции. Z-score: $z = (x - \mu) / (\sigma + \varepsilon)$, где $\varepsilon = 10^{-8}$ защищает от деления на ноль.

\textbf{Этап 4. Скользящее среднее.} Вычисляется через \texttt{np.cumsum} за $O(N)$. Разница соседних значений темпа --- через \texttt{np.diff} с последующим \texttt{np.insert}.

\textbf{Этап 5. Feature engineering.} $\texttt{speed\_kmh} = \texttt{dist} / (\texttt{pace}/60)$, $\texttt{cal\_per\_min} = \texttt{cal} / \texttt{pace}$. Защита от деления на ноль через \texttt{np.maximum}. Нефинитные значения заменяются медианой.

\textbf{Этап 6. Условная агрегация.} Маска: зона $\geq 4$ AND дистанция $>$ медиане AND темп $<$ медиане. Для каждого спортсмена считаются mean, median, p90 по \texttt{cal\_per\_min} при выполнении условия. Результат сохраняется в два CSV.

\textbf{Этап 7. Лаговые признаки.} \texttt{np.roll(arr, 1)} сдвигает массив на 1 позицию (последний становится первым). Первый элемент заменяется на самого себя. Вычисляется разность \texttt{pace - pace\_lag1}.

\textbf{Этап 8. Робастная замена выбросов.} Для каждой группы вычисляются Q1, Q3, IQR. Точки вне $[Q1 - 1.5 \cdot IQR,\ Q3 + 1.5 \cdot IQR]$ заменяются медианой группы. Итог: доля замен и статистика по группам.

\textbf{Этап 9. Логическая целостность.} Правила: зона 5 при pace $> 8.5$ --- невозможна, зона 1 при pace $< 3.2$ --- невозможна. Выявленные противоречия корректируются (зона устанавливается в 3).

\textbf{Этап 10. Сжатие редких категорий.} Категории zone с частотой $< 1\%$ помечаются как 0 (прочие). Сохраняется финальный массив \texttt{final\_preprocessed.npy}.
\end{flushleft}

\newpage
\begin{flushleft}
\section*{Выходные файлы}

\begin{center}
\begin{tabular}{|l|l|}
\hline
\textbf{Файл} & \textbf{Содержимое} \\
\hline
\texttt{stage1\_structured.npy}        & Исходные данные в структурированном dtype \\
\texttt{stage3\_normalized.npy}        & Данные + z\_dist после нормализации \\
\texttt{conditional\_stats.csv}        & mean, median по cal\_per\_min при условии \\
\texttt{conditional\_stats\_with\_p90.csv} & То же + p90 \\
\texttt{final\_preprocessed.npy}       & Финальный массив (все 12 полей) \\
\hline
\end{tabular}
\end{center}

\section*{Вывод программы (пример)}

\begin{verbatim}
Строк: 500000
Память: 7000000 байт (6.676 Мбайт)
NaN по полям: {'dist': 0, 'pace': 0, 'cal': 0}
Inf по полям: {'dist': 0, 'pace': 0, 'cal': 0}
Аномалий: 312 (0.06%)
Групп: 50
Размеры групп: min=9850, max=10200
Доля роста pace: 49.87%, падения: 49.91%, без изменений: 0.22%
Распределение sign: {-1: 249553, 0: 1112, 1: 249335}
Доля замен после IQR: 0.82%
Замены по группам (первые 10): [(1, 82), (2, 78), ...]
Логических нарушений: 0 (0.00%)
Редкие категории zone: []
Готово. Сохранены файлы: stage1_structured.npy, ...
\end{verbatim}
\end{flushleft}

\begin{flushleft}
\newpage
\section*{Вывод}
\noindent
В ходе лабораторной работы я освоил полный цикл предобработки данных с помощью NumPy. Реализованы десять этапов: загрузка в структурированный массив с явным dtype, векторная очистка аномалий, z-нормализация внутри групп спортсменов, скользящее среднее через cumsum, feature engineering (скорость, калории/мин), условная агрегация с сохранением в CSV, лаговые признаки через np.roll, робастная IQR-замена выбросов, проверка логической согласованности зон пульса и темпа, сжатие редких категорий. Все операции реализованы через векторные функции NumPy без явных Python-циклов там, где это возможно.
\end{flushleft}

\end{document}
